\subsection{Confirmatory study: forecasting versus tuned smoothing}\label{sec:v10}
The three analyses in Section~\ref{sec:posthoc} were performed after the test
results were seen, and the EWMA coefficient 0.1 was itself chosen after v8, so
on their own they cannot settle the question. We therefore pre-registered a
third study (v10, \code{docs/service_v10_protocol.md}, written before the first
cell was run) to test the remaining question directly: once \emph{both} quality
sources are smoothed with the coefficient that is best for each of them, does
the GNN forecast still raise PDR?

Every arm uses the route-pinning rule, the binary, the flags and the locked GNN
weights of v8, with stationary relays. The EWMA family smooths the current
measurement, $q_t=\alpha d_t+(1-\alpha)q_{t-1}$; the GNN family is given the
same knob on the forecast, $q_t=\beta\hat g_t+(1-\beta)q_{t-1}$, where
$\beta=1$ is exactly \code{gnn-route} of v8. The coefficient of each family was
selected on a separate tuning phase of 160 episodes whose seeds are audited
disjoint from every earlier plan including v8x, under a single pre-registered
criterion: the highest n5/n7-balanced mean gradual PDR within the family, ties
broken towards the larger coefficient. Tuning returned 160/160 valid cells and
selected $\alpha^\ast=\beta^\ast=0.1$. Validation returned 60/60 valid cells and
passed its technical gate, and the test returned 144/144 valid cells with route
integrity, forecast traces and relay stationarity verified in every episode.

\begin{table}[htbp]
\centering\small
\begin{tabularx}{\linewidth}{@{}Yp{3.4cm}p{1.9cm}p{2.4cm}@{}}
\toprule
\textbf{Comparison (v10 test)}&\textbf{PDR difference (pp), 95\% CI}&\textbf{n5 / n7}&\textbf{Verdict}\\\midrule
H1: \code{gnn-tuned} $-$ \code{ewma-tuned}&$+0.08$ [$-0.11$, 0.29]&0.18 / $-0.01$&No difference detected\\
H2: \code{ewma-tuned} $-$ \code{lq-route}&$+0.90$ [0.51, 1.35]&1.19 / 0.61&Descriptive only\\
\code{gnn-tuned} $-$ \code{lq-route}&$+0.98$ [0.60, 1.42]&--&Descriptive only\\
\bottomrule
\end{tabularx}
\caption{The v10 test: 40 paired gradual blocks, 10{,}000 block bootstrap resamples, balanced across n5 and n7.}\label{tab:v10test}
\end{table}

\begin{table}[htbp]
\centering\small
\begin{tabularx}{\linewidth}{@{}Yp{2.7cm}p{2.1cm}p{2.0cm}p{1.5cm}@{}}
\toprule
\textbf{Arm}&\textbf{PDR (\%, mean $\pm$ SD)}&\textbf{Route changes (total)}&\textbf{Disconnected (s)}&\textbf{5~s MSE}\\\midrule
\code{lq-route} (EWMA 0.5)&95.08 $\pm$ 2.72&385&31.18&0.00129\\
\code{ewma-tuned} (EWMA 0.1)&95.98 $\pm$ 2.44&83&30.27&0.00166\\
\code{gnn-tuned} (GNN, $\beta=0.1$)&96.07 $\pm$ 2.36&61&30.52&0.00123\\
\bottomrule
\end{tabularx}
\caption{Descriptive statistics for the v10 gradual test; the five-second MSE is computed on the same \code{lq-route} traces, balanced across n5 and n7.}\label{tab:v10means}
\end{table}

The primary endpoint detected no difference (Table~\ref{tab:v10test}): once both
sides are tuned, the GNN leads the EWMA by 0.08 percentage points with a 95\%
confidence interval of $[-0.11, 0.29]$, and it is higher in 22 of 40 gradual
blocks, lower in 16 and tied in 2. Because H1 reached no directional conclusion,
H2 is reported descriptively even though its interval lies above zero; all three
arms deliver the same PDR on the static blocks. Between the two tuned arms the
difference in disconnected time is $+0.25$ seconds, CI $[-0.53, 1.00]$, and the
background-flow PDR difference is $-0.08$ percentage points, CI $[-0.22, 0.06]$.

This confirms the exploratory conclusion through a pre-registered test: the
$+0.66$ percentage point gain of v8 does not require a learned forecaster,
because a tuned one-parameter EWMA reaches the same level ($+0.90$ against the
same untuned EWMA 0.5 baseline). The direction of forecast error also repeats at
confirmatory scale: on the same test traces EWMA 0.1 has the worst five-second
MSE of the three arms (0.00166 against 0.00123 for the smoothed GNN) yet no
worse PDR (Table~\ref{tab:v10means}). Route-change counts fall monotonically
with the amount of smoothing, from 385 to 83 and 61, consistent with the
proposed mechanism: the route-pinning controller behaves better on a less noisy
input, and a learned forecaster smooths as a side effect. The interval
$[-0.11, 0.29]$ still admits a small difference in either direction, so v10 does
not establish equivalence; it shows that whatever gain is attributable to
forecast information as such is smaller than 0.3 percentage points in this
scenario family.
